\section*{الفصل \ref{ch:b2:fragment-orbitals} --- مدارات الشظايا والجزيئات المتعددة الذرات}

\begin{solution}{exo:b2:fragment-orbitals:1}
$h_1 + h_2$ و$h_1 - h_2$. المستوي $xz$، الذي ينصّف الزاوية ويكون
عموديًا على الجزيء، يبادل بين الهيدروجينين: فيغيّر $h_1 - h_2$
إشارته، ولا يغيّرها $h_1 + h_2$.
\end{solution}

\begin{solution}{exo:b2:fragment-orbitals:2}
ستة مدارات تكافؤ (2s وثلاثة 2p للأكسجين، و1s للهيدروجينين)، وثمانية إلكترونات
تكافؤ، وأربعة مدارات ممتلئة.
\end{solution}

\begin{solution}{exo:b2:fragment-orbitals:3}
مستويات التكافؤ الممتلئة مدار $a_1$ واحد ومجموعة من ثلاثة مدارات $t_2$
متساوية الطاقة: طاقتان، فحزمتان. وحزمة $t_2$ تنزع إلكترونات من ثلاثة
مدارات تحمل ستة إلكترونات، مقابل مدار واحد وإلكترونين للمدار $a_1$.
\end{solution}

\begin{solution}{exo:b2:fragment-orbitals:4}
تتطلب رابطة $\pi$ أن يكون مدارا p للكربونين متوازيين، وهذا يضع
مستويي \ce{CH2} كليهما في مستوٍ واحد.
\end{solution}

\begin{solution}{exo:b2:fragment-orbitals:5}
تفقد الأمونيا إلكترونًا من زوجها الحر، وهو \omterm{def:b2:diatomic-mos:bonding}{مدار غير رابط} في معظمه على
النيتروجين؛ ويفقده الميثان من مدار $t_2$ رابط. والإلكترون غير الرابط يقع أعلى،
أقرب إلى المستوى الذري، من الإلكترون الرابط: فيُنزع بسهولة أكبر.
\end{solution}

\begin{solution}{exo:b2:fragment-orbitals:6}
للماء الخطي مدارا p ممتلئان غير رابطين متساويا الطاقة؛ والانحناء يُخفض أحدهما
ويرفع مستوى رابطًا قليلًا: ربح صافٍ مع ثمانية إلكترونات. أما \ce{BeH2}
فليس له إلا أربعة، تملأ \omterm{def:b2:diatomic-mos:bonding}{المدارين الرابطين}؛ والمدار الذي كان الانحناء سيخفضه
فارغ، بينما تفقد المستويات الرابطة من تداخلها: فالشكل الخطي هو الأفضل.
\end{solution}

\begin{solution}{exo:b2:fragment-orbitals:7}
كل رابطة C--H أو C--C على كربون مجاور للمركز الكاتيوني يمكن أن تصطف مع
المدار p الفارغ وتعطي استقرارًا بنحو $2\beta^2/\Delta$. وليس للكاتيون \ce{CH3+} أي
جار من هذا النوع؛ أما كاتيونات الإيثيل والإيزوبروبيل وثالثي البوتيل فلها مجموعة ميثيل واحدة واثنتان وثلاث
مجاورة للمركز، تقدم كل منها رابطة مصطفة: فيزداد الاستقرار بهذا
الترتيب.
\end{solution}

\begin{solution}{exo:b2:fragment-orbitals:8}
من أجل مستويين متساويين يكون الانشطار $2|\beta|$ ويكسب إلكترونا $\pi$ مقدار $2|\beta|$،
المتناسب مع $\cos\theta$. وينخفض إلى النصف عند $\theta = 60^\circ$ وينعدم عند
$90^\circ$.
\end{solution}

\begin{solution}{exo:b2:fragment-orbitals:9}
للبوران ستة إلكترونات تكافؤ: تملأ \omterm{def:b2:diatomic-mos:bonding}{المدارات الرابطة} الثلاثة، ويبقى المدار p
للبور العمودي على المستوي فارغًا؛ فلا ربح من
التهرّم. وللأمونيا إلكترونان زائدان: يشغلان \omterm{def:b2:diatomic-mos:bonding}{المدار غير الرابط} في
الشكل المستوي، وكما في الماء، يسمح التهرّم له بالامتزاج مع المدار 2s 
والتركيب الهيدروجيني فينخفض.
\end{solution}

\begin{solution}{exo:b2:fragment-orbitals:10}
تعطي المعادلات القرنية من أجل $c_1 = c_2 = c_3$ القيمة $\alpha + 2\beta$؛ ويعطي التركيبان
المتعامدان معه $\alpha - \beta$ (متساويين في الطاقة). إلكترونان في $\alpha + 2\beta$: الطاقة
$2\alpha + 4\beta$، أدنى من $\ce{H2} + \ce{H+}$ ($2\alpha + 2\beta$) لأن $\beta < 0$.
\end{solution}

\begin{solution}{exo:b2:fragment-orbitals:11}
$\pi_1$: المعاملات الثلاثة كلها بإشارة واحدة، بلا عقدة (رابط)؛ $\pi_2$: إشارتان متعاكستان
على الكربونين الطرفيين وعقدة في الوسط (غير رابط)؛ $\pi_3$: إشارات
متناوبة، وعقدتان (مضاد للربط). وللأنيون أربعة إلكترونات $\pi$: فيكون $\pi_2$
أعلى مستوياته المشغولة، وكثافته على الكربونين الطرفيين وحدهما.
\end{solution}

\begin{solution}{exo:b2:fragment-orbitals:12}
في كل من المستويين اللذين يحويان المحور، تعطي مدارات p الثلاثة المتوازية (O، C، O)
\omterm{def:b2:diatomic-mos:bonding}{مدارًا رابطًا} \omterm{def:b2:diatomic-mos:bonding}{ومدارًا غير رابط} (على الأكسجينين وحدهما) \omterm{def:b2:diatomic-mos:bonding}{ومدارًا مضادًا للربط}: ستة
مدارات $\pi$ في المجموع. ويمتلئ \omterm{def:b2:diatomic-mos:bonding}{المداران الرابطان} \omterm{def:b2:diatomic-mos:bonding}{والمداران غير الرابطين}: ثمانية
إلكترونات $\pi$ (رابطتا $\pi$ وزوجان حران للأكسجين في صورة لويس).
\end{solution}

\begin{solution}{pb:b2:fragment-orbitals:1}
\textbf{1.} ثمانية: ستة من الأكسجين، وواحد من كل هيدروجين.
\textbf{2.} $h_1 + h_2$ و$h_1 - h_2$.
\textbf{3.} 2s و$2\mathrm p_z$.
\textbf{4.} $2\mathrm p_y$.
\textbf{5.} $2\mathrm p_x$، العمودي على المستوي الجزيئي.
\textbf{6.} ستة؛ أربعة منها ممتلئة.
\textbf{7.} $h_1 - h_2$، الذي يغيّر إشارته، مثل $2\mathrm p_z$، عبر المستوي
العمودي على المحور عند الأكسجين.
\textbf{8.} $h_1 + h_2$.
\textbf{9.} $2\mathrm p_x$ و$2\mathrm p_y$، العموديان على المحور: زوج غير رابط
متساوي الطاقة.
\textbf{10.} (2s، $h_1 + h_2$) رابط$^2$، و($2\mathrm p_z$، $h_1 - h_2$) رابط$^2$، ثم
$2\mathrm p_x^2\,2\mathrm p_y^2$.
\textbf{11.} لا: الزوج المتساوي الطاقة ممتلئ، وكل إلكترون مقترن.
\textbf{12.} عند طاقة مدار 2p للأكسجين: فهو غير رابط.
\textbf{13.} $2\mathrm p_y$ (محور الجزيء الخطي على طول $z$، والانحناء في $yz$).
\textbf{14.} يمتزج مع 2s و$h_1 + h_2$ فينخفض: المستوى $3a_1$.
\textbf{15.} المستوى الرابط ($2\mathrm p_z$، $h_1 - h_2$)، الذي يفقد شيئًا من التداخل حين
يغادر الهيدروجينان المحور: فيصير $1b_2$.
\textbf{16.} المستوى الممتلئ الذي ينخفض يكسب أكثر مما يخسر الآخر: مع ثمانية
إلكترونات يكون الشكل المنحني هو الأكثر استقرارًا.
\textbf{17.} $104.48^\circ$، بين $90^\circ$ (روابط من مدارات p وحدها) و$109.5^\circ$:
فالمدار 2s يسهم بقسط في الارتباط.
\textbf{18.} $2\mathrm p_x$، الذي يبقى \omterm{def:b2:diatomic-mos:bonding}{المدار غير الرابط} $1b_1$.
\textbf{19.} أربع.
\textbf{20.} من $1b_1$، المدار $2\mathrm p_x$ غير الرابط للأكسجين.
\textbf{21.} نزع إلكترون غير رابط لا يكاد يغيّر الروابط: فللأيون
هندسة الجزيء ولا يُثار إلا قليل من الطاقة الاهتزازية. أما نزع إلكترون
رابط فيغيّر أطوال الروابط وزواياها، فتمتد الحزمة على
مستويات اهتزازية كثيرة.
\textbf{22.} \qty{12.62}{eV} مقابل \qty{13.62}{eV}: يحمل الأكسجين في الماء شحنة
سالبة جزئية مسحوبة من الهيدروجينين، ترفع مستوياته 2p.
\textbf{23.} ليسا مستويين متساويين: فتركيباهما هما المداران $3a_1$ 
و$1b_1$، المختلفا الطاقة. والصورتان تصفان الكثافة الإلكترونية الكلية نفسها؛
لكن المدارات غير الموضعية وحدها تعطي طاقات التأين.
\textbf{24.} للماء \textbf{أربع} حزم تأين تكافؤية، أدناها طاقةً عند
$\boldsymbol{\approx \qty{12.6}{eV}}$.
\end{solution}
